\section{Summary}\label{sec:summary}

\paragraph{What was done.}
The transcript analysis rests on two language models that rate passages of machine transcripts on a bank of
questions, each with five described levels from 0 to 4, and on eight dimensions computed from those ratings. In
two studies people rated the same text the models rated: the first 250 words of a 500-word passage, shown as
plain text.
\begin{itemize}
\item \textbf{The first study} took twelve rating questions, the highest-loading questions of the eight
  dimensions. Each of 160 excerpts was rated on each question by five people: 9{,}648 ratings by 205 raters. The
  question is the unit of its verdicts.
\item \textbf{The second study} took sixteen questions, two markers for each dimension, and made the dimension
  the unit. 211 raters gave 13{,}248 ratings of 160 excerpts, two passages from each of 80 randomly drawn
  recordings. Every rater read a different mix of excerpts, so that a rater's general level can be told apart
  from the excerpts and taken out, and half the raters were shown no sample answers.
\end{itemize}
In both studies the design, the checks that decide whose ratings count, the estimators and the rules for the
verdicts were fixed before collection. A question or dimension is \verd{usable} when the mean of five raters per
question has a reliability of 0.60 or more and at least one model's score correlates with that mean at 0.50 or
more.

\paragraph{First study: results fixed in advance.}
\begin{itemize}
\item Four questions are usable: positive tone, disappointment, self-critical wording and self-disclosure. The
  mean of five raters has a reliability of 0.76 to 0.85 and the models correlate with it at 0.72 to 0.89
  (Figure~\ref{fig:verdicts}). They are the marker questions of dimensions F2 and F3.
\item No question is \verd{misread}: wherever raters agree with each other, the models follow them.
\item Eight questions are not settled. What holds them back is how far raters agree with each other; the better
  model correlates with the mean of five at 0.53 to 0.71 on every one.
\item The two models follow people equally well: 0.673 for \jev{} and 0.670 for \drex{}.
\end{itemize}

\paragraph{Second study: results fixed in advance.}
\begin{itemize}
\item \textbf{Four dimensions are usable, with the whole interval above the lines}: F1 (direct, organized
  argument), F2 (positive tone versus disappointment), F3 (self-disclosure and self-criticism) and F6 (energetic
  and angry versus calm). The composite of ten raters has a reliability of 0.71 to 0.86 and the better model's
  dimension score correlates with it at 0.72 to 0.82 (Figure~\ref{fig:s2verdicts}).
\item \textbf{Three dimensions lie at the lines}: F4 (care and regard for others), F5 (narrative versus
  rebuttal) and F7 (humor and imaginative language).
\item \textbf{One dimension has an unclear rubric}: F8 (hedging and open-mindedness versus certainty). Raters do
  not agree at all on one of its two questions.
\item No dimension is misread, and the two models are again level.
\item Showing raters sample answers after the practice excerpts does not pull them towards the models.
\end{itemize}

\paragraph{Added after the results had been read,} with less weight.
\begin{itemize}
\item \textbf{The first study repeats.} Across the twelve questions rated in both studies, how well a model
  follows people correlates at 0.95 and 0.90 between the studies. With both studies pooled, calm joins the four
  usable questions.
\item \textbf{The settled verdicts are firm.} They hold when the choice of raters is counted as well as the
  choice of recordings, when the reliability is computed without a statistical model, and when the ratings are
  treated as ordered categories.
\item \textbf{A model is closer to people's shared reading than one person is.} Against what an unlimited panel
  of readers would say, the models correlate at 0.81 to 0.86 on average; a single rating at 0.51 to 0.56.
\item \textbf{The models follow the passage, not only the recording.} The difference between a model's scores
  of two passages of one recording follows the difference between people's ratings of them.
\item \textbf{Each dimension score follows its own dimension} more closely than any other, clearly so for six
  of the eight.
\item \textbf{People see the differences the models see} between videos and livestreams and over the years,
  with one exception: at the same score for care and regard for others, people rate livestreams half a standard
  deviation lower than videos.
\item \textbf{The models' levels are not people's levels.} One model rates hedging about a scale point higher
  than people; the other rates most questions lower.
\item \textbf{A third model does the same.} d1.2-ultra, released after both studies, was put to the same test
  on the same excerpts. It is level with the two models on the eight dimensions, slightly ahead on single
  questions, shares their misses, and its levels are closer to people's (Section~\ref{sec:third}).
\end{itemize}

\paragraph{Together.} For four of the eight dimensions the models' rankings of excerpts are supported by
people's ratings; for three the evidence lies at the lines of the rule; for one the pair of questions does not
work in people. Rankings are supported; levels are not (Table~\ref{tab:standing}).

\paragraph{How to read this report.}
Section~\ref{sec:question} says what is tested and by which rule. Sections~\ref{sec:first} and~\ref{sec:second}
give the two studies, Section~\ref{sec:further} the analyses added afterwards, Section~\ref{sec:third} a third
model, Section~\ref{sec:discussion} what
the studies support and what they do not, and Section~\ref{sec:materials} the participants, the data and the
code. Intervals in brackets are 95\% intervals from resampling excerpts or recordings unless stated otherwise.
